\begin{abstract}
Se documenta el diseño, la implementación y los ensayos de vuelo de un controlador de actitud y altura para un cuadricóptero F450 (1.0022\,kg) gobernado por un ESP32-S3. El controlador sigue la formulación del curso: linealización por realimentación en cascada (ángulo $\to$ velocidad angular $\to$ torque), control de altura con cancelación de gravedad y mezclador por inversión de la matriz de rotores, y se extiende con un integrador, un super-twisting opcional y una identificación recursiva de la autoridad de control. Los parámetros del modelo se obtuvieron experimentalmente: inercias con péndulo bifilar ($I_x=9.71$, $I_y=10.56$, $I_z=20.08\times10^{-3}$\,kg\,m$^2$), el coeficiente de empuje con un banco de motor y el de par a partir de la relación empuje/par de la hélice. Se analizan dos vuelos de la versión V\_11 y se los compara con la V\_10. En el vuelo B el dron despegó y se sostuvo en el aire con $|\phi|<10^{\circ}$ el 78\,\% del tiempo y un pico de $57^{\circ}$ recuperado sin desarme; el estimador de altura mejoró frente a V\_10 (desviación respecto del ultrasonido $0.12$\,m frente a $0.26$--$1.39$\,m), aunque diverge junto al suelo cuando se queda sin medida válida. Se discuten las limitaciones: margen de empuje reducido (relación empuje/peso 1.43), discrepancia entre la curva de empuje de banco y el vuelo, y ausencia de medición de batería y de posición lateral.
\end{abstract}

\begin{IEEEkeywords}
cuadricóptero, linealización por realimentación, control en cascada, identificación de parámetros, filtro de Kalman, super-twisting, ESP32-S3.
\end{IEEEkeywords}
